\documentclass[../main]{subfiles}
\begin{document}
\newpage

\subsubsection{\texttt{oe} --- Output Error}

\begin{lstlisting}
theta, m = pysib.oe(u, y, nb, nf, nz)
\end{lstlisting}

Identifies an OE model by minimizing the prediction-error criterion.
The model structure is

\[
  y(t) = \frac{B(q^{-1})}{F(q^{-1})}\,u(t - n_k) + e(t),
\]

where $e(t)$ is white noise.
For this structure the one-step-ahead predictor equals the
noise-free simulated output, so the criterion is

\[
  V(\theta) = \frac{1}{N}\sum_t
  \left[y(t) - \frac{B(q^{-1})}{F(q^{-1})}\,u(t-n_k)\right]^2.
\]

This is nonlinear in $\theta$ because $F$ appears in the denominator.
An ARX estimate provides the initial guess; the optimizer then
alternates steepest-descent and Newton steps (implemented in a C
extension, with LAPACK used for the Newton step).

\subsubsection*{Arguments}
\begin{center}
\begin{tabular}{lll}
  \toprule
  Name & Type & Description \\
  \midrule
  \texttt{u}  & array\_like & Input signal. \\
  \texttt{y}  & array\_like & Output signal. \\
  \texttt{nb} & int & Number of $B$ parameters. \\
  \texttt{nf} & int & Number of $F$ parameters. \\
  \texttt{nz} & int & Input delay $n_k \geq 0$. \\
  \bottomrule
\end{tabular}
\end{center}

\subsubsection*{Returns}
\begin{center}
\begin{tabular}{lll}
  \toprule
  Name & Type & Value \\
  \midrule
  \texttt{theta} & ndarray & $[b_0,\ldots,b_{n_b-1},\; f_1,\ldots,f_{n_f}]$ \\
  \texttt{m}     & dict    & \texttt{B}, \texttt{F} free; \texttt{A}$=$\texttt{C}$=$\texttt{D}$=[1]$ \\
  \bottomrule
\end{tabular}
\end{center}

\end{document}
